\begin{table*}[t]
\centering
\caption{Taxonomy of sustainable marine biofouling-control strategies. NR denotes not reported in the available corpus.}
\label{tab:taxonomy}
\small
\begin{tabularx}{\textwidth}{l l l X X}
\toprule
Control mode & Representative systems & Target stage & Main mechanism & Principal limitation \\
\midrule
Antifouling & Metal oxides, biocides, bio-derived compounds & Settlement and early biofilm & Toxicity, reactive species, deterrence, or chemistry-mediated inhibition & Release and non-target effects; efficacy is condition-dependent \\
Fouling-release & Silicone, silicone composites, low-modulus surfaces & Established soft and hard fouling & Low adhesion and interfacial fracture under motion or cleaning & Requires retained smoothness, elasticity, and suitable operating conditions \\
Anti-adhesion & Hydrated, zwitterionic, heparin, hydrogel, amphiphilic surfaces & Bacteria, diatoms, larvae & Hydration layer, steric barrier, or reduced interfacial interaction & Ageing, contamination, and mechanical durability \\
Self-cleaning & Photocatalytic, UV, ciliary, lubricant-infused surfaces & Deposits and early fouling & Stimulus-driven degradation, movement, or low-shear release & Energy, optical access, lubricant loss, or scale-up constraints \\
Active maintenance & Grooming, waterjet, ultrasonic, robot, laser & Established fouling and niche areas & Mechanical or energy-assisted removal & Coating wear, waste capture, access, and operational cost \\
\bottomrule
\end{tabularx}
\end{table*}
